\subsubsection{Kịch bản 1: GPU trong mạng nội bộ}
\label{sec:scenario1}

\textbf{Bối cảnh áp dụng.} Bệnh viện đã có hạ tầng đầy đủ, kể cả GPU đặt tại chỗ, và
có đội ngũ công nghệ thông tin vận hành được một cụm máy chủ. Mọi xử lý, kể cả suy luận
AI, gán nhãn và huấn luyện, diễn ra trong mạng nội bộ; đám mây chỉ dùng để sao lưu và
lưu trữ dài hạn. Kiến trúc của kịch bản này là Hình~\ref{fig:arch-s1}.

\textbf{Điểm khác biệt.} Không có điểm biến thiên nào được bật: dữ liệu không rời
bệnh viện cho mục đích AI, nên luồng suy luận không cần khử định danh hay mã hoá riêng
và không cần cơ chế đặc biệt để tối ưu độ trễ; GPU đọc thẳng bản làm việc trên cổng
trung gian qua mạng nội bộ. Cổng trung gian vẫn được giữ, dù về lý thuyết dịch vụ suy
luận có thể đọc thẳng từ kho đối tượng dùng chung, vì ba lý do: PACS cần một đích
C-STORE để đẩy dữ liệu sang mà không cấp quyền truy cập cho AI; tầng nghiên cứu cần một
điểm khử định danh duy nhất; và giữ cổng giúp Kịch bản 1 và Kịch bản 2 dùng cùng quy tắc
xuất dữ liệu và cùng giao diện suy luận, nên một bệnh viện có thể chuyển kịch bản mà
không cấu hình lại PACS.

\textbf{Triển khai.} Các thành phần tầng lâm sàng và tầng dữ liệu chạy trên cụm máy chủ
của bệnh viện với dự phòng N+1 (Bảng~\ref{tab:ha}). Các nút GPU suy luận và GPU huấn luyện
tách riêng: suy luận phải chạy liên tục để phục vụ cấp cứu, còn huấn luyện là các tác vụ
nặng chạy theo đợt, nên dùng chung GPU sẽ khiến một lần huấn luyện làm chậm cảnh báo cho
ca cấp cứu. Vùng mạng của GPU không có đường ra Internet. Số lượng và cấu hình GPU phụ
thuộc thời gian suy luận thực tế của mô hình trên một thể tích 3D và lưu lượng ca chụp
của bệnh viện, nên cần được đo trên mô hình đã chọn trước khi chốt.

\textbf{Luồng 2: suy luận AI.}
\begin{description}[leftmargin=1.2cm,style=nextline]
\item[2.1] Quy tắc xuất trong dcm4chee chọn các chuỗi ảnh AI cần (CT trục lát mỏng) và
      đẩy sang cổng trung gian bằng C-STORE. Nếu cổng không sẵn sàng, tác vụ được xếp hàng
      và thử lại. Khi bác sĩ cần so sánh, ca cũ của cùng bệnh nhân cũng có thể được đẩy
      sang theo quy tắc xuất ca cũ hoặc theo yêu cầu.
\item[2.2] Cổng lưu chuỗi ảnh (chỉ mục trong PostgreSQL, tệp trên SeaweedFS). Khi chuỗi
      không nhận thêm ảnh nào trong một khoảng thời gian cấu hình được, cổng phát sự kiện
      ``đã đủ ảnh''; bộ điều phối tạo một tác vụ suy luận, gán hạn chót và độ ưu tiên, và
      bỏ qua tác vụ trùng nếu một nút cổng khác đã tạo cho cùng ca.
\item[2.3] Dịch vụ suy luận đọc thể tích từ cổng, tiền xử lý theo đúng quy trình đã dùng
      khi huấn luyện, chạy mô hình đã duyệt trên GPU và tạo mặt nạ vùng xuất huyết. Cờ ưu
      tiên được suy ra trước từ mặt nạ; báo cáo định lượng được tạo sau.
\item[2.4] Mặt nạ (DICOM SEG) và báo cáo (DICOM SR) kèm đủ siêu dữ liệu bắt buộc
      (Mục~\ref{sec:aicontract}) trở về cổng; cổng gửi chúng vào PACS bằng C-STORE tới AE
      \texttt{AI\_RESULTS}, vốn chỉ cho phép lưu. Kết quả nằm cùng ca chụp gốc trên PACS và
      bác sĩ thấy lớp phủ trên trình xem như một kết quả chưa xác nhận.
\item[2.5] Nếu nghi có xuất huyết, hệ thống gửi cảnh báo tới bác sĩ qua kênh trong mạng
      nội bộ và gửi cờ ưu tiên tới HIS/RIS (HL7 ORU hoặc FHIR) để ca được đưa lên đầu
      worklist. Ca không được gắn cờ giữ nguyên vị trí trong worklist.
\end{description}

\textbf{Khi có sự cố.} Khi GPU, dịch vụ suy luận hoặc cổng ngừng hoạt động, ca chụp
vẫn được lưu và đọc bình thường, chỉ thiếu cảnh báo và lớp phủ; trình xem và worklist
hiển thị trạng thái ``AI tạm ngừng'' để không ai hiểu nhầm việc không có cờ là âm tính.
Khi cổng hoạt động trở lại, PACS tự gửi lại các ca đang xếp hàng. Một tác vụ vượt quá hạn
chót (Bảng~\ref{tab:latency}) được đánh dấu ``không có kết quả AI'' thay vì chờ vô hạn.

\textbf{Vấn đề còn mở.} (i) Giá trị cụ thể của khoảng chờ ``đã đủ ảnh'' và chu kỳ quét
tác vụ của PACS để đạt mục tiêu 90 giây cho ca cấp cứu, cần đo trên máy chụp thật. (ii)
Mọi máy CT của bệnh viện có gửi chuỗi tái tạo lát mỏng về PACS hay không; nếu không, mô
hình 3D không có đầu vào phù hợp. (iii) HIS/RIS hiện tại có nhận và hiển thị cờ ưu tiên
qua HL7 ORU hoặc FHIR hay không.
